\documentclass[11pt,a4paper]{article}

\usepackage[utf8]{inputenc}
\usepackage[T1]{fontenc}
\usepackage[round,sort,comma]{natbib}
\usepackage{amsmath,amssymb}
\usepackage{booktabs}
\usepackage{listings}
\usepackage{xcolor}
\usepackage{hyperref}
\usepackage{microtype}
\usepackage{geometry}
\geometry{margin=1in}

\lstset{
  basicstyle=\ttfamily\small,
  frame=single,
  breaklines=true,
  keepspaces=true,
}

\setlength{\parindent}{0pt}
\setlength{\parskip}{0.5em}

\title{Program Generation via Reinforcement Learning\\with Enumeration Warm-Start}
\author{Yiding Song}
\date{\today}

\begin{document}
\maketitle

\section*{Problem Statement}

We are interested in the following problem. Suppose we have access to:
\begin{itemize}
  \item A set $F$ of functions $f_1, f_2, \ldots, f_n$ with types
        $\tau_1 \to \upsilon_1, \tau_2 \to \upsilon_2, \ldots, \tau_n \to \upsilon_n$,
        where we can query the function implementation (i.e.\ $f_i(x)$ can be computed
        for arbitrary well-typed inputs $x$) but not necessarily the function description
        (e.g.\ we do not know the mathematical or programmatic form of the functions).
  \item A reward function $R: \mathrm{program} \to \mathbb{R}$ that maps a program to a
        real-valued reward, where a program is a well-typed composition of functions from
        $F$ (e.g.\ this can reward a function for desirable properties such as being
        non-constant).
\end{itemize}

{\it How do we synthesise a large number of programs that are diverse and achieve high rewards?}

This document describes a two-phase program synthesis solution to the problem: bottom-up enumeration
followed by reinforcement learning.

\section*{Overview}

The pipeline has two phases:

\begin{enumerate}
  \item \textbf{Bottom-up enumeration} --- exhaust all programs up to a certain size,
        quotienting out semantic duplicates (as measured by function behaviour on a test
        suite), creating a corpus of behaviourally diverse programs.
  \item \textbf{RL exploration} --- a neural policy learns from the corpus (behavioural
        cloning warm-start), then generates new programs via policy gradient, using a
        priority queue buffer of the best programs found so far
        \citep{abolafia2018priority}.
\end{enumerate}

\section{The Program Grammar}

Programs are expressions in a simply-typed lambda calculus extended with conditionals and
a fixed set of primitive functions. The expression language is:
\[
  e \;::=\; c \;\mid\; x \;\mid\; f(e_1, \ldots, e_k)
         \;\mid\; \lambda(x_1, \ldots, x_k).\, e
         \;\mid\; \mathrm{if}\ e_1\ \mathrm{then}\ e_2\ \mathrm{else}\ e_3
\]
where $c$ ranges over literals (integers, booleans, empty lists), $x$ over variables, and
$f$ over the primitive functions in $F$.

\subsection{Types}

The base types are $\mathrm{int}$ and $\mathrm{bool}$. From these, list types
$\mathrm{list}[\sigma]$ and function types $\sigma \to \tau$ are formed. In the current
instantiation the concrete types in use are:
\[
  \mathrm{int},\quad \mathrm{bool},\quad \mathrm{list}[\mathrm{int}],\quad
  \mathrm{list}[\mathrm{bool}],\quad \mathrm{list}[\mathrm{list}[\mathrm{int}]]
\]
All programs synthesised by the pipeline have type
$\mathrm{list}[\mathrm{int}] \to \mathrm{list}[\mathrm{int}]$.

\subsection{Polymorphic Functions}

Some primitives in $F$ are \textbf{polymorphic}: their types contain type variables that
can be instantiated to any concrete type. For example:
\[
  \mathrm{map} : (T_1 \to T_2) \to \mathrm{list}[T_1] \to \mathrm{list}[T_2]
\]
Here $T_1$ and $T_2$ are type variables. A \textbf{type substitution} $\theta$ maps each
type variable to a concrete type; applying $\theta$ to the signature of $\mathrm{map}$
yields a monomorphic instantiation. For instance,
$\theta = \{T_1 \mapsto \mathrm{int},\, T_2 \mapsto \mathrm{bool}\}$ gives:
\[
  \mathrm{map}_\theta : (\mathrm{int} \to \mathrm{bool})
    \to \mathrm{list}[\mathrm{int}] \to \mathrm{list}[\mathrm{bool}]
\]
The set of valid substitutions for a function is determined by the concrete types in use.
Each instantiation is treated as a separate monomorphic function during enumeration.

\subsection{Higher-Order Functions}

A function is \textbf{higher-order} if one or more of its arguments has a function type.
\texttt{map} is higher-order: its first argument is a function $T_1 \to T_2$. Such
arguments cannot be looked up directly in the program bank (which stores complete
expressions, not open function values); instead they must be synthesised as lambda
expressions, as described in the Higher-Order Applications section below.

\section{Part 1: Bottom-Up Enumeration}

\subsection{Size Metric}

Every AST node $e$ is assigned an integer size $\lvert e \rvert$ inductively:

\begin{center}
\begin{tabular}{ll}
\toprule
Node & Size \\
\midrule
Literal, variable, empty list
  & $1$ \\
Application $(f\ a_1\ \cdots\ a_k)$
  & $1 + \lvert a_1 \rvert + \cdots + \lvert a_k \rvert$ \\
Lambda $(\lambda\ (p_1\ \cdots\ p_k)\ e)$
  & $1 + \lvert e \rvert$ \\
Conditional $(\mathrm{if}\ c\ t\ e)$
  & $1 + \lvert c \rvert + \lvert t \rvert + \lvert e \rvert$ \\
\bottomrule
\end{tabular}
\end{center}

Enumeration proceeds by increasing size: programs of size $s$ are built exclusively from
sub-expressions of size less than $s$.

\subsection{Observational Equivalence and the Program Bank}

We remove program duplicates by comparing their behaviour on a test suite. Fix a test suite $T = \{t_1, \ldots, t_m\}$ of $m$
concrete inputs. The \textit{fingerprint} of a closed program $p$ of type
$\tau \to \upsilon$ is:
\[
  \Phi(p) = \bigl(p(t_1),\ \ldots,\ p(t_m)\bigr) \in (\upsilon \cup \{\dagger\})^m
\]
where $\dagger$ records a failed evaluation (runtime error or type error). Two programs
$p, q : \tau \to \upsilon$ are observationally equivalent if $\Phi(p) = \Phi(q)$.

\paragraph{Open terms and lambda arguments.}
However, not all programs are closed. To enumerate an expression like
$\mathrm{map}\ (\lambda y.\ e)\ x$, the enumerator must synthesise candidate bodies $e$.
These bodies are \emph{open terms}: they may reference $y$ (the lambda parameter) freely,
alongside the outer $x$. Storing them in the bank allows the same body to be reused
across different higher-order functions, avoiding redundant work.

In general, for an open term $e$ with free variables
$y_1 : \sigma_1, \ldots, y_k : \sigma_k$ (beyond the primary input $x$), the fingerprint
is computed over $T \times V_1 \times \cdots \times V_k$, where each $V_i$ is a finite
set of probe values for type $\sigma_i$:
\[
  \Phi(e)
    = \bigl(e(t,\, v_1,\ldots,v_k)\bigr)_{(t,\, v_1,\ldots,v_k)
        \,\in\, T \times V_1 \times \cdots \times V_k}
\]
Two open terms are considered observationally equivalent iff they agree on this full
Cartesian product.

The \textbf{program bank} $\mathcal{B}$ maps each pair $(\tau, s)$ to a set of programs,
containing exactly one representative per observational equivalence class:
\[
  \mathcal{B}[\tau][s]
    = \bigl\{[p]_\sim \;\big|\; \vdash p : \tau,\; \lvert p \rvert = s\bigr\}
\]

\paragraph{Default test suite} ($m = 10$):
\begin{lstlisting}
[], [0], [3,1,2], [1,1,1,1], [5,4,3,2,1],
[1,2,3,4,5,6,7,8], [10,-3,7,7,0], [2,8,3,8,2,3],
[0,1,0,1,0], [42]
\end{lstlisting}

\subsection{Base Case}

The bank is seeded with size-1 expressions:
\begin{itemize}
  \item \textbf{Integer constants} $c \in C_{\text{seed}}$ of type $\mathrm{int}$
  \item \textbf{Boolean constants} $\top, \bot$ of type $\mathrm{bool}$
  \item \textbf{Empty lists} $[\,]$ of type $\mathrm{list}[\sigma]$ for each list type $\sigma$
  \item \textbf{Input variable} $x : \mathrm{list}[\mathrm{int}]$ (since all of our
        programs are of type $\mathrm{list}[\mathrm{int}] \to \mathrm{list}[\mathrm{int}]$,
        this is the variable that will be bound to the input list in a program like
        \texttt{length x})
\end{itemize}

\subsection{First-Order Applications}

For a first-order function $f : A_1 \times \cdots \times A_k \to R$ and target size $s$,
the arguments must satisfy $\sum_{i=1}^k s_i = s - 1$ with each $s_i \geq 1$. For every
such composition:
\[
  \mathcal{B}[R][s]
    \;\mathrel{+}=\;
    \bigl\{f(a_1, \ldots, a_k) \;\big|\; a_i \in \mathcal{B}[A_i][s_i]\bigr\}
    \!\Big/\!\sim
\]
Polymorphic functions are pre-expanded into all valid monomorphic instantiations (as
described in the Grammar section); each instantiation is then enumerated independently
using the formula above.

\subsection{Higher-Order Applications and Nested Lambdas}

When an argument position requires a function type $A_1 \times \cdots \times A_k \to B$,
the lambda argument is synthesised inline rather than looked up in the flat bank. Given a size
budget $b$ for the lambda, an extended typing context
$\Gamma' = \Gamma \cup \{y_i : A_i\}_{i=1}^k$ is formed, and a child bank $\mathcal{B}'$
is enumerated recursively over $\Gamma'$ up to size $b - 1$. The lambda candidates are
then:
\[
  \bigl\{(\lambda\ (y_1 \cdots y_k)\ e) \;\big|\; e \in \mathcal{B}'[B][b-1]\bigr\}
\]
Child banks are keyed by $(\Gamma', b)$ and cached so that the same lambda type appearing
as an argument to several functions at the same size is only built once.

\paragraph{Nesting depth.}
$\ell$ counts the number of enclosing lambda scopes. Higher-order calls within a child
context are only attempted when $\ell \leq \ell_{\max}$ (default $\ell_{\max} = 2$),
preventing combinatorial explosion from arbitrarily deep lambda nesting while still
enabling patterns with three levels of nesting such as
$\mathrm{map}\ (\lambda y.\ \mathrm{map}\ (\lambda z.\ \mathrm{map}\ (\lambda w.\ \cdots)\ z)\ y)\ x$.

\subsection{Quality Filtering}

After enumeration, the corpus is obtained by retaining programs that pass three
conditions. Let $\mathcal{F}(p) \subseteq \{1,\ldots,m\}$ denote the indices on which
$p$ does not crash. A program $p$ is retained iff:
\begin{enumerate}
  \item $|\mathcal{F}(p)| \geq 3$ \quad (non-crashing on sufficiently many inputs)
  \item $|\{\,p(t_i) \mid i \in \mathcal{F}(p)\}| \geq 2$ \quad (non-constant output)
  \item $\dfrac{|\{\,p(t_i) \mid i \in \mathcal{F}(p)\}|}{|\mathcal{F}(p)|}
        \geq \nu_{\min}$ \quad (output variability, default $\nu_{\min} = 0.3$)
\end{enumerate}

\subsection{Key Parameters}

\begin{center}
\begin{tabular}{lll}
\toprule
Parameter & Default & Role \\
\midrule
$s_{\max}$        & 5         & Maximum AST size enumerated \\
$\ell_{\max}$     & 2         & Maximum lambda nesting depth for HOF arguments \\
$C_{\text{seed}}$ & $\{0,1,\ldots,10\}$
                             & Integer literals available as base-case expressions \\
$\nu_{\min}$      & 0.3       & Corpus variability threshold \\
$T$               & 10 inputs & Test suite for fingerprinting and deduplication \\
\bottomrule
\end{tabular}
\end{center}

\section{Part 2: Reinforcement Learning}

The RL phase trains a neural policy on the enumerated corpus to generate good programs
recursively from the root, making one decision per AST node. This is framed as a Markov
Decision Process (MDP) over a grammar-constrained action space.

\subsection{MDP Formulation}

\paragraph{State.}
A synthesis state $s$ is a tuple $(\tau, \Gamma, f_{\text{par}}, i, \delta, \ell)$:

\begin{center}
\begin{tabular}{ll}
\toprule
Component & Description \\
\midrule
$\tau$                  & Required type at the current node \\
$\Gamma$                & Typing context: variables in scope with their types \\
$f_{\text{par}},\ i$   & Parent function and argument index (contextual features) \\
$\delta$                & Remaining depth budget (decremented at each recursive call) \\
$\ell$                  & Current lambda nesting depth \\
\bottomrule
\end{tabular}
\end{center}

\paragraph{Actions.}
The valid action set $\mathcal{A}(s)$ depends on $s$:

\begin{center}
\begin{tabular}{ll}
\toprule
Action & Condition \\
\midrule
Literal integer $c$  & $\tau = \mathrm{int}$ \\
Literal boolean      & $\tau = \mathrm{bool}$ \\
Empty list           & $\tau$ is a list type \\
Variable $y$         & $y \in \mathrm{dom}(\Gamma)$ with $\Gamma(y) = \tau$ \\
Apply $f$ (instance) & $\delta > 0$; return type of $f$ matches $\tau$;
                       if $f$ is higher-order, $\ell < \ell_{\max}^{\mathrm{RL}} = 2$ \\
Lambda               & $\tau$ is a function type \\
If                   & $\delta > 0$ \\
\bottomrule
\end{tabular}
\end{center}

\paragraph{Episode.}
Generation begins from the state
$s_0 = (\mathrm{list}[\mathrm{int}] \to \mathrm{list}[\mathrm{int}],\
\varnothing,\ \cdot,\ \cdot,\ \delta_{\max},\ 0)$
and proceeds by recursively expanding each chosen action into child states:
\begin{itemize}
  \item $\mathrm{Apply}\ f(A_1, \ldots, A_k) \to R$: recurse on
        $(\tau \leftarrow A_i,\ \delta - 1)$ for each $i$
  \item $\mathrm{Lambda}\ (\tau_1 \times \cdots \times \tau_k \to \sigma)$: recurse on
        $(\tau \leftarrow \sigma,\ \Gamma \cup \{y_i : \tau_i\},\ \delta - 1,\ \ell + 1)$
  \item $\mathrm{If}$: recurse on $(\mathrm{bool},\ \delta-1)$ then twice on
        $(\tau,\ \delta-1)$
\end{itemize}
If any recursive call encounters an empty action set, the episode fails and returns no
program.

\subsection{Reward}

Given a completed program $p$, let $\Phi(p)$ denote its fingerprint on $T$, and let
$\mathcal{K}$ be the set of fingerprints already known (from the enumeration corpus and
prior RL discoveries). The reward is:
\[
  R(p) = \mathrm{var}(\Phi(p)) \times
    \begin{cases}
      1.0 & \text{if } \Phi(p) \notin \mathcal{K} \\
      0.1 & \text{otherwise}
    \end{cases}
\]
where $\mathrm{var}(\Phi(p)) = |\{p(t_i)\}| / m$ is the fraction of distinct outputs.
Programs with $|\mathcal{F}(p)| < 3$ or $|\{p(t_i)\}| < 2$ receive $R(p) = 0$ and are
discarded.

The novelty multiplier incentivises exploration of behaviours not yet covered, while the
$0.1$ floor ensures the policy is not penalised for rediscovering known programs.

\subsection{Priority Queue Buffer}

The buffer follows the priority queue training scheme of \citet{abolafia2018priority}.
$\mathcal{B}_{\mathrm{buf}}$ maintains the top-$K$ programs by reward (default
$K = 5000$). It is a bounded min-heap: a new program $p$ is admitted only if
$R(p) > R_{\min}(\mathcal{B}_{\mathrm{buf}})$, evicting the current minimum. Duplicate
fingerprints are excluded.

Training samples are drawn uniformly from $\mathcal{B}_{\mathrm{buf}}$ with no recency or
priority weighting.

\subsection{Policy Network}

The policy $\pi_\theta$ maps a state $s$ to a distribution over actions via a masked
softmax. A network $g_\theta$ produces a logit for each action; invalid actions are
excluded by setting their logits to $-\infty$ before normalising:
\[
  \pi_\theta(a \mid s)
    = \frac{\exp\bigl(g_\theta(s)_a\bigr)}
           {\displaystyle\sum_{a' \in \mathcal{A}(s)} \exp\bigl(g_\theta(s)_{a'}\bigr)},
  \qquad a \in \mathcal{A}(s)
\]
where $g_\theta : \mathcal{S} \to \mathbb{R}^{|\mathcal{A}|}$ is a two-layer feedforward
network:
\[
  g_\theta(s) = W_2\, \mathrm{ReLU}(W_1\, \mathrm{enc}(s) + b_1) + b_2
\]

\paragraph{State encoder.}
Six features of $s$ are each embedded or projected into $\mathbb{R}^d$, then concatenated
and projected back to $\mathbb{R}^d$:

\begin{center}
\begin{tabular}{ll}
\toprule
Feature & Encoding \\
\midrule
$\tau$              & Learned embedding over a type vocabulary \\
$f_{\mathrm{par}}$  & Learned embedding over the function vocabulary \\
$i$ (arg index)     & Learned embedding, clamped to $[0,7]$ \\
$\delta$            & Learned embedding, clamped to $[0,15]$ \\
$\ell$              & Learned embedding, clamped to $[0,3]$ \\
$\Gamma$            & Linear projection of a per-type variable-count vector \\
\bottomrule
\end{tabular}
\end{center}

Default $d = 64$, hidden dimension 128. Invalid actions are masked to $-\infty$ before
softmax; the valid set $\mathcal{A}(s)$ is cached by $(\tau, \Gamma, \delta > 0, \ell)$.

\subsection{Training Phase 1: Behavioural Cloning Warm-Start}

Each corpus program $p$ induces a sequence of state--action pairs $\{(s_t, a_t)\}$ by a
DFS traversal of its AST, recording at each node the state that a top-down policy would
have been in and the action it would have taken. Aggregating over all corpus programs
gives a dataset $\mathcal{D}$. The policy is pre-trained by minimising:
\[
  \mathcal{L}_{\mathrm{BC}}(\theta)
    = -\mathbb{E}_{(s,\, a)\,\sim\,\mathcal{D}}\bigl[\log \pi_\theta(a \mid s)\bigr]
\]
(50 epochs, batch size 64, Adam at $\eta = 10^{-3}$). Transitions involving actions
outside the vocabulary are dropped.

\subsection{Training Phase 2: RL Loop}

Each iteration alternates between a sampling phase and a gradient phase.

\paragraph{Sampling} (32 episodes per iteration): run an episode under $\pi_\theta$;
compute $\Phi(p)$ and $R(p)$; if $R(p) > 0$, attempt insertion into
$\mathcal{B}_{\mathrm{buf}}$; if $\Phi(p)$ is novel, add it to $\mathcal{K}$.

\paragraph{Gradient update} (8 steps per iteration): sample a batch from
$\mathcal{B}_{\mathrm{buf}}$ and minimise the reward-weighted log-likelihood loss
\citep{abolafia2018priority}:
\[
  \mathcal{L}_{\mathrm{RL}}(\theta)
    = -\mathbb{E}_{p\,\sim\,\mathcal{B}_{\mathrm{buf}}}
        \left[ R(p) \sum_{t} \log \pi_\theta(a_t \mid s_t) \right]
\]
No baseline is subtracted. Gradient norms are clipped to 1.0 before the Adam step
($\eta = 10^{-4}$).

\subsection{Trajectory Extraction}

The dataset $\mathcal{D}$ (used in both BC and RL) is constructed by a DFS that is the
inverse of the episode procedure. Given a complete AST, each node yields a state--action
pair:
\begin{itemize}
  \item $\lambda$-node $\to$ action \textbf{Lambda}, recurse into body at $\ell + 1$
  \item Application $f(a_1, \ldots, a_k)$ $\to$ action \textbf{Apply} $f$, recurse into
        each $a_i$ with $f_{\mathrm{par}} = f$, $i$ set accordingly
  \item If-node $\to$ action \textbf{If}, recurse into condition then branches
  \item Leaf $\to$ the appropriate literal or variable action
\end{itemize}
For polymorphic functions, the type instantiation is recovered by matching the inferred
return type of $f$ against $\tau$ in the current state.

% \section{Limitations}

% \paragraph{Partial application / currying.} All functions are treated as taking a fixed number of arguments. The value $(= y) : \mathrm{int} \to \mathrm{bool}$ cannot be produced directly; a lambda wrapper $(\lambda z.\ z = y)$ is required instead.

% \paragraph{Lambda in function position.} The Apply action always names a grammar function; applying an arbitrary expression in the function position (e.g., $(\lambda y.\ \cdots)\ \mathrm{arg}$) is not supported by either the enumerator or the MDP. This pattern, equivalent to a let-binding, is uncommon in practice.

% \paragraph{Constant coverage.} Integer literals are restricted to $C_{\mathrm{seed}}$. Programs requiring values outside this set cannot be generated without explicitly extending $C_{\mathrm{seed}}$.

% \paragraph{Size and depth limits.} Enumeration halts at $s_{\max}$ and generation at depth $\delta_{\max}$. These are configuration parameters rather than architectural constraints, but complex programs (e.g., nested folds) may exceed them at default settings.

\bibliographystyle{plainnat}
\bibliography{references}

\end{document}
